% ===================================================================
%  جدول نمبر 12 — اسٹیشن۔ --- ریلیں۔ --- جہاز۔
%
%  Traduction, faite en 2026, en ourdou d'aujourd'hui, sur l'ido de
%  texto/io. Regles de la colonne : voir 10-jadval-01.tex. Une seule
%  numerotation. Ecrit avec outils/ordo.py sous les yeux. Le « Noto »
%  d'ouverture appartient a l'apparat et n'a pas de cle propre,
%  comme du cote ido ; la note de bas de page du feuillet 85 en a
%  une, t12-noto-1, et le bloc t12-14-1 se poursuit apres elle sous
%  la MEME cle.
%
%  قلی EST LE DEUXIEME MOT DU RELEVE QUI A VOYAGE DANS L'AUTRE SENS.
%  Le tableau 5 avait montre برآمدہ, la veranda : l'anglais l'a pris
%  a l'ourdou et non l'inverse, si bien que la colonne « emprunte »
%  un mot qui est chez lui. Le porteur (16) fait le second : قلی est
%  ourdou et hindi, l'anglais en a fait « coolie » et l'a exporte
%  dans le monde entier. Une colonne qui compte ses emprunts doit
%  compter aussi ceux-la, et dans le bon sens.
%
%  LE CHEMIN DE FER EST LE CONTRAIRE EXACT DE LA MARINE DU TABLEAU
%  10. Sur mer, tout ce qui precedait la vapeur avait un nom persan
%  et tout ce qu'elle a invente etait anglais. Ici la vapeur EST
%  l'objet, et pourtant l'ourdou nomme presque tout : ڈبّا le wagon
%  — le mot ordinaire pour une boite —, پٹری le rail, پھاٹک la
%  barriere, اوقات نامہ l'indicateur, مال گاڑی le train de
%  marchandises, بھاپ la vapeur, سیٹی le sifflet, پہیے les roues,
%  انتظار گاہ la salle d'attente, بیت الخلا les latrines, منشی le
%  commis. Restent انجن, پلیٹ فارم, ٹکٹ, گارڈ, سگنل, بوائلر, پسٹن :
%  sept mots. La difference avec la marine tient a une chose et une
%  seule — le chemin de fer indien a ete construit tot et a servi
%  tout le monde, pendant que la marine a vapeur restait le metier
%  d'un corps d'officiers qui parlait anglais. Le mot suit l'USAGE,
%  pas la technique.
%
%  ET LES SEPT EMPRUNTS SE TIENNENT PAR LA MAIN. انجن, پلیٹ فارم,
%  ٹکٹ, گارڈ, سگنل sont les cinq mots qu'un voyageur lit ECRIT sur
%  place, en anglais, depuis 1861 ; بوائلر et پسٹن sont les deux
%  pieces qu'un voyageur ne voit jamais. Aucun des sept ne designe
%  une chose qu'on manie sans savoir lire. C'est une regle qu'on
%  pourra verifier ailleurs : l'emprunt s'installe la ou la chose
%  porte une etiquette.
%
%  TRENTE PAIRES, ET UNE PARADE QUI N'EST PAS LA RELATIVE. Au bloc
%  t12-06-1 l'ido pose la sortie (32) avant la valise (24) — « li
%  hastis al ekireyo kun sua valizo en manuo ». L'ourdou met le
%  complement de maniere avant le verbe, donc la valise remonterait
%  devant la sortie, et « جو » ne peut rien ici : ce n'est pas une
%  relation de tete a complement, c'est un ordre de phrase. La parade
%  est le rejet apres tiret, qui est du bon ourdou ecrit et que la
%  colonne emploie depuis le tableau 7. Deux outils, deux emplois :
%  la relative pour la tete, le tiret pour la circonstance.
% ===================================================================

%%K t12-apar-1 apar
\VUsaut{4.0mm}
\VUtitre{15.0pt}{60}{جدول نمبر 12}
\VUsaut{2.4mm}
\VUpk{11.4pt}{40}{اسٹیشن۔ --- ریلیں۔ --- جہاز۔}
\VUsaut{3.4mm}
\textsc{نوٹ۔} --- \textit{ہر ملک میں مستعمل اوقات کے نقشے مختلف ہیں،}
\textit{اِس لیے ہم مثال کے طور پر} \VUgras{فرانسیسی نقشے} \textit{کے
اوقات لیتے ہیں۔}

%%K t12-01-1 p
1. --- میں ابھی جرمنی کے سفر سے ہو کر آیا ہوں۔ روانگی سے پہلے میں نے
ایک \VUgras{اوقات نامہ} \textsuperscript{(15)} خریدا تاکہ گاڑیوں کے
چلنے کا وقت معلوم ہو۔ یہ دیکھ لینے کے بعد کہ سب سے موزوں گاڑی پیرس سے
رات نو بجے چلتی ہے اور اسٹراسبورگ ایک بج کر تیرہ منٹ پر پہنچتی ہے،
میں نے اپنا سفر تیار کیا، اور اگلے دن خود کو اسٹیشن پہنچوا دیا۔

%%K t12-02-1 p
2. --- جب میں بڑے روانگی ہال میں داخل ہوا، ایک بوڑھے دیہاتی
\VUgras{پادری} \textsuperscript{(2)} کے پیچھے، جو ہاتھ میں اپنا
\VUgras{سفری تھیلا} \textsuperscript{(3)} اٹھائے ہوئے تھا، بہت سے
\VUgras{مسافر} \textsuperscript{(1)} پہلے ہی آ چکے تھے۔

%%K t12-02-2 p
چونکہ میں ایک \VUgras{تار} \textsuperscript{(5)} بھیجنا چاہتا تھا،
میں نے ایک \VUgras{نگران} \textsuperscript{(6)} سے
\VUgras{تار گھر} \textsuperscript{(4)} کا پتا پوچھا۔

%%K t12-03-1 p
3. --- \VUgras{انتظار گاہ} \textsuperscript{(7)} میں جانے سے پہلے میں
واپسی کا \VUgras{ٹکٹ} \textsuperscript{(9)} لینے گیا۔

%%K t12-04-1 p
4. --- \VUgras{کھڑکی} \textsuperscript{(8)} پر مجھے زیادہ دیر انتظار
نہ کرنا پڑا، کیونکہ وہاں لوگ کم تھے۔ ایک \VUgras{سپاہی}
\textsuperscript{(10)}، جو چھٹی پر جا رہا تھا، مہربانی سے مجھے اپنی
باری دے گیا ؛ چنانچہ مجھے اتنا وقت مل گیا کہ \VUgras{اسٹیشن ماسٹر}
\textsuperscript{(13)} سے کچھ باتیں پوچھ لوں، جو \VUgras{نگرانی کے
افسر} \textsuperscript{(12)} اور ڈیوٹی پر موجود \VUgras{فوجی پولیس کے
سپاہی} \textsuperscript{(11)} سے باتیں کر رہے تھے۔

%%K t12-tit-1 sub
\VUpk{10.2pt}{40}{سامان کی درج کاری۔}
\VUsaut{3.0mm}

%%K t12-05-1 p
5. --- \VUgras{کتابوں کے کھوکھے} \textsuperscript{(14)} سے ایک اخبار
خرید کر میں نے اپنا \VUgras{سامان} \textsuperscript{(17)} تلوایا، جو
ایک \VUgras{قلی} \textsuperscript{(16)} ابھی ابھی \VUgras{سامان کی
ریڑھی} \textsuperscript{(19)} پر لایا تھا ؛ \VUgras{وہ منشی}
\textsuperscript{(22)}، جو \VUgras{تولنے کے آلے} \textsuperscript{(21)}
پر مامور ہے، نے مجھے بتایا کہ میرا صندوق پینتیس کلو کا ہے ؛ یعنی پانچ
کلو \VUgras{زائد وزن}، جس کے لیے مجھے \VUgras{سامان کی کھڑکی}
\textsuperscript{(23)} پر اضافی رقم دینی پڑی۔

%%K t12-06-1 p
6. --- چونکہ میں وقت سے بہت پہلے پہنچ گیا تھا، مجھے فرصت تھی کہ
اسٹیشن پر مسافروں کی آمد و رفت دیکھوں۔ کچھ لوگ اپنا معمولی سامان
\VUgras{امانتی سامان گھر} \textsuperscript{(25)} میں رکھوا رہے تھے یا
واپس لے رہے تھے ؛ کچھ اور \VUgras{اوقات کے نقشے}
\textsuperscript{(26)} دیکھ رہے تھے۔ آنے والے مسافر \VUgras{نکاس}
\textsuperscript{(32)} کی طرف لپک رہے تھے — ہاتھ میں اپنا
\VUgras{اٹیچی} \textsuperscript{(24)} لیے۔ کچھ لوگ \VUgras{سامان
بانٹنے کی جگہ} \textsuperscript{(27)} پر انتظار کر رہے تھے اور اپنی
\VUgras{رسید} \textsuperscript{(29)} دکھا کر اپنے صندوق،
\VUgras{بکس} \textsuperscript{(28)} یا \VUgras{ٹوکرے}
\textsuperscript{(30)} واپس لے رہے تھے۔

%%K t12-07-1 p
7. --- \VUgras{محصول کے اہلکار} \textsuperscript{(33)} گٹھڑیوں کو
غور سے دیکھ رہے تھے کہ کہیں کوئی ظاہر کرنے کی چیز تو نہیں۔

%%K t12-08-1 p
8. --- کچھ مسافر \VUgras{کھانے کے کمرے} \textsuperscript{(34)} کی طرف
گئے، جہاں ایک \VUgras{بیرا} \textsuperscript{(35)} اُنھیں پیش کرنے میں
مستعد تھا۔ اُنھیں جلدی کرنی ہے، کیونکہ \VUgras{ریل گاڑی}
\textsuperscript{(36)}، جو تیز رفتار ہے، اسٹیشن پر زیادہ دیر نہیں
ٹھہرے گی۔

%%K t12-09-1 p
9. --- پھر میں روانگی کے پلیٹ فارم پر گیا اور سیدھا \VUgras{ڈبّا}
\textsuperscript{(37)} ڈھونڈا تاکہ سفر کے دوران ڈبّا بدلنے پر مجبور نہ
ہوں۔ میں ایک راہ دار والی گاڑی میں چڑھا، جس میں تمباکو نوشوں کے لیے ایک
\VUgras{خانہ} \textsuperscript{(38)} ہے اور ایک خانہ « اکیلی خواتین »
کے لیے مخصوص۔

%%K t12-tit-2 sub
\VUpk{10.2pt}{40}{رات میں روانگی۔}
\VUsaut{3.0mm}

%%K t12-10-1 p
10. --- \VUgras{پائیدان} \textsuperscript{(40)} پر چڑھ کر،
\VUgras{دروازہ} \textsuperscript{(39)} بند کر کے، \VUgras{پردہ}
\textsuperscript{(41)} لپیٹ کر اور اپنا چھوٹا سامان \VUgras{جالی}
\textsuperscript{(43)} پر رکھ کر میں \VUgras{بنچ}
\textsuperscript{(42)} پر بیٹھ گیا۔ \VUgras{لیمپ والے}
\textsuperscript{(44)} کو لیمپ جلاتے دیکھ کر میں نے سوچا کہ مجھے رات
ڈبّے ہی میں گزارنی ہو گی، اور میں نے \VUgras{تکیوں}
\textsuperscript{(45)} کے کرائے پر مامور اہلکار کو بلا کر ایک تکیہ
مانگا۔

%%K t12-11-1 p
11. --- ٹکٹ چیکر ٹکٹ دیکھنے آیا۔ میں نے اُس سے پوچھا کہ ہم اب تک
روانہ کیوں نہیں ہوئے، جب کہ وقت ہو چکا ہے۔ اُس نے جواب دیا کہ
\VUgras{گارڈ} \textsuperscript{(46)} \VUgras{سامان کے ڈبّے}
\textsuperscript{(47)} میں سائیکلیں رکھوانے میں لگا ہے۔ اِس کے علاوہ
سارے \VUgras{پاؤں گرم رکھنے والے} \textsuperscript{(48)} ابھی بدلے
نہیں گئے۔

%%K t12-12-1 p
12. --- مگر ہم جلد ہی روانہ ہونے والے تھے، کیونکہ \VUgras{کوئلہ
جھونکنے والا} \textsuperscript{(50)} اپنی \VUgras{کوئلہ گاڑی}
\textsuperscript{(49)} پر سے کوئلہ اٹھا رہا تھا تاکہ \VUgras{انجن}
\textsuperscript{(54)} کی آگ بھڑکائے، اور \VUgras{ڈرائیور}
\textsuperscript{(51)} صرف \VUgras{اسٹیشن کے نائب افسر}
\textsuperscript{(52)} کے اشارے کا منتظر تھا، جو \VUgras{پلیٹ فارم}
\textsuperscript{(53)} پر کھڑا تھا۔

%%K t12-13-1 p
13. --- انجن کا \VUgras{بوائلر} \textsuperscript{(56)} دبی دبی گرج رہا
تھا ؛ \VUgras{چمنی} \textsuperscript{(55)} \VUgras{بھاپ}
\textsuperscript{(60)} میں ملے دھوئیں کے دریا اُگل رہی تھی۔ ایک
چیختی \VUgras{سیٹی} \textsuperscript{(59)} گونجی اور \VUgras{پسٹن}
\textsuperscript{(58)} حرکت میں آئے، اُس بھاری مشین کے \VUgras{پہیوں}
\textsuperscript{(57)} کو دھکیلتے ہوئے۔

%%K t12-tit-3 sub
\VUsaut{3.0mm}
\VUpk{10.2pt}{40}{روانگی!}
\VUsaut{3.0mm}

%%K t12-14-1 p
14. --- \VUgras{پٹریوں} \textsuperscript{(64)} پر دوڑتی گاڑی کی رفتار
تیز ہوتی گئی ؛ \VUgras{پلیٹ فارم} \textsuperscript{(65)} غائب ہو گئے ؛
ہم \VUgras{بیت الخلا} \textsuperscript{(66)} کے پاس سے گزرے، اور اُن
ڈبّوں کے پاس سے بھی جو دوسری \VUgras{لائن} \textsuperscript{(63)} پر
کھڑے تھے : \VUgras{سونے کے ڈبّے} \textsuperscript{(67)}،
\VUgras{کھانے کا ڈبّا} \textsuperscript{(68)}، \VUgras{ڈاک کا ڈبّا}
\textsuperscript{(69)}۔

%%K t12-noto-1 noto
\VUnotes{143.2mm}{%
\textit{(*) نوٹ۔} --- \textit{ظاہر ہے کہ سفر کا یہ بیان جنگ سے پہلے
کی سرحد کے مطابق ہے۔}%
}

%%K t12-15-1 p
15. --- پھر \VUgras{مال گودی} \textsuperscript{(71)} ہماری آنکھوں کے
سامنے سے گزری۔ شیڈوں کے نیچے اہلکار وہ \VUgras{مال}
\textsuperscript{(72)} لاد رہے تھے جو سست رفتار گاڑی سے بھیجا جاتا ہے۔
\VUgras{مزدوروں کے دستے} \textsuperscript{(76)} \VUgras{پانی کی ٹنکی}
\textsuperscript{(73)} کے پاس \VUgras{جانوروں کے ڈبّے}
\textsuperscript{(75)} اور \VUgras{کھلے ڈبّے} \textsuperscript{(79)}
اِدھر اُدھر کر رہے تھے تاکہ ایک \VUgras{مال گاڑی}
\textsuperscript{(80)} بنائیں۔

%%K t12-16-1 p
16. --- ہم آخری \VUgras{پٹری کا کانٹا} \textsuperscript{(78)} پار کر
گئے، جس کے پاس \VUgras{کانٹے والا} کھڑا تھا ؛ ہم \VUgras{سگنل}
\textsuperscript{(86)} سے آگے نکل گئے اور اسٹیشن دور جا کر غائب ہو
گیا۔

%%K t12-17-1 p
17. --- جلد ہی ہم ایک \VUgras{لوہے کے پُل} \textsuperscript{(74)} سے
گزرے، جو \VUgras{دریا} \textsuperscript{(81)} کے اوپر سے جاتا ہے۔ وہ
پُل پار کرنے کے بعد ہم دریا کے ساتھ ساتھ چلے، جس پر زور دار
\VUgras{ٹگ} \textsuperscript{(82)} بھاری \VUgras{مال بردار کشتیاں}
\textsuperscript{(83)} کھینچ رہے تھے۔ ہم نے ایک \VUgras{مسافر جہاز}
\textsuperscript{(84)} بھی دیکھا، جب وہ ایک \VUgras{تیرتے گھاٹ}
\textsuperscript{(85)} سے لگ رہا تھا۔

%%K t12-18-1 p
18. --- کئی اسٹیشنوں کے پاس سے تیزی سے گزرنے کے بعد ہم چند
\VUgras{سرنگوں} \textsuperscript{(87)} سے گزرے اور \VUgras{پھاٹک
والوں} کے بے شمار \VUgras{گھروندے} \textsuperscript{(88)} پیچھے چھوڑ
آئے۔ گاڑی کے جھولے میں جھولتے ہوئے مجھے گہری نیند آ گئی۔

%%K t12-tit-4 sub
\VUsaut{3.0mm}
\VUpk{10.2pt}{40}{ڈوئچ-آورِیکور اسٹیشن۔}
\VUsaut{3.0mm}

%%K t12-19-1 p
19. --- مجھے اچانک ایک اہلکار کی آواز نے جگا دیا، جو پکار رہا تھا :
« ڈوئچ-آورِیکور! \textsuperscript{(*)} سب اتر جائیں! » محصول کا معائنہ
ہوا اور سرحد کی وہ ساری اکتا دینے والی کارروائیاں۔ مجھے اپنی جیبی
گھڑی اسٹیشن کی اُس \VUgras{بڑی گھڑی} \textsuperscript{(89)} کے مطابق
کرنی پڑی، جو پچپن منٹ آگے تھی ؛ ابھی ابھی جاگا ہوا میں بمشکل
\VUgras{بڑی سوئی} \textsuperscript{(90)} کو \VUgras{چھوٹی}
\textsuperscript{(91)} سے الگ پہچان سکا ؛ اور \VUgras{گھڑی کا چہرہ}
\textsuperscript{(92)} خود صبح کی دھند سے بالکل دھندلا تھا۔

%%K t12-20-1 p
20. --- جب تک محصول کے اہلکار معائنہ کرتے رہے، میں جماہیاں لیتے ہوئے
وہ \VUgras{اشتہار} \textsuperscript{(93)} پڑھتا رہا جو اسٹیشن کی
دیواروں کو سجاتے ہیں۔ وقت گزارنے کے لیے میں نے ناشتہ کیا، جرمن
\VUgras{ریل کے جال} \textsuperscript{(94)} کا نقشہ دیکھا تاکہ معلوم ہو
کہ اسٹراسبورگ ابھی کتنا دور ہے، اور اپنے ڈبّے میں واپس آ گیا — اِس
اُمید سے تازہ دم کہ اپنے سفر کی منزل جلد پہنچ جاؤں گا۔

\VUsaut{6.0mm}
\VUfilet{22mm}
